
%==========================================================================
\section{Conclusion et travaux futurs}\label{sec:conclusion}
%==========================================================================

Nous avons implémenté eALS~\cite{he2016eals} en PySpark - une version RDD fondée sur
des broadcasts, adossée à une référence NumPy séquentielle - et vérifié que les deux
produisent le même modèle à $\CorrWorst$ près, quel que soit le nombre de blocs, ce qui
est le pendant empirique de l'affirmation de l'article selon laquelle eALS est
trivialement parallélisable \emph{sans approximation}. En plus de cela, nous avons
implémenté l'évaluation leave-one-out sur catalogue entier et une campagne
expérimentale dont la sortie brute se trouve dans \texttt{results/*.csv}.

Les résultats que nous considérons comme la substance de ce rapport :

\begin{enumerate}\itemsep2pt
\item L'affirmation de complexité tient, dans le régime pour lequel elle est
énoncée. Sur Amazon Movies, le coût d'une itération entre $K=32$ et $K=128$ croît de
$\GrowthEalsAmz\times$ pour eALS et $\GrowthAlsAmz\times$ pour l'ALS implicite de
MLlib, sur les mêmes données, la même machine et le même nombre de cœurs. eALS passe
de $\RatioAtMinAmz\times$ plus lent à $K=\KAtRatioMinAmz$ à $\RatioAtMaxAmz\times$
plus rapide à $K=\KAtRatioAmz$. Ajuster les deux formes de complexité avec des
constantes libres reproduit les mesures à $R^2>0,999$, mais place le croisement entre
le terme linéaire et le terme quadratique d'eALS à $K\approx\FitCrossEalsAmz$ plutôt
qu'au $\CrossKAmz$ prédit par le décompte de flops : le terme creux est limité par la
mémoire et le terme dense relève de BLAS, et un décompte de flops ne peut pas voir la
différence.
\item La scalabilité forte est réelle mais modeste : $\SpeedupBig\times$ sur 8 cœurs
pour ml-10m, cohérent avec une fraction série d'Amdahl de $\SerialBig$. La limite
n'est pas l'algorithme - le parallélisme est exact et le travail se découpe
proprement - mais la bande passante mémoire partagée, la désérialisation des
broadcasts dans chaque worker Python, et la surcharge Python par tâche.
\item La pondération sensible à la popularité fonctionne : à $\alpha=0,4$, HR@100
passe de $\HRu$ à $\HRe$ sur Amazon Movies, soit $\PopGain$\,\% de mieux que le
réglage uniforme $\alpha=0$, toutes choses égales par ailleurs.
\item Sur la métrique qui compte en pratique - la précision pour un budget de temps
donné - eALS atteint en $\EalsTimeToAls$\,s la meilleure précision que l'ALS de
MLlib n'atteint jamais (HR@100 $=\AlsBestHR$ après $\AlsTimeToBest$\,s), puis continue
de s'améliorer.
\end{enumerate}

\paragraph{Travaux futurs.} Quatre directions, dans l'ordre où nous les prendrions.
\emph{(i) Métriques de couverture et de biais de popularité.} L'ajout le plus utile,
de loin : mesurer la couverture du catalogue et la concentration de Gini en fonction
de $\alpha$, et voir quelle part du gain de HR relève réellement d'une meilleure
modélisation des préférences.
\emph{(ii) Une implémentation co-partitionnée.} Remplacer le broadcast de $Q$ par un
échange par blocs partitionnés où chaque ligne de facteur n'est envoyée qu'aux blocs
qui la référencent - le schéma qu'utilise l'ALS de MLlib - et ré-exécuter l'étude
de scalabilité sur un vrai cluster, où le coût du broadcast cesse d'être une copie
mémoire.
\emph{(iii) Références BPR et RCD}, pour tester l'asymétrie rappel/précision que nous
ne pouvons actuellement que citer depuis l'article.
\emph{(iv) Float32 et une boucle interne Numba ou Cython.} Le noyau est limité par la
mémoire ; diviser par deux la largeur de chaque gather devrait rapporter près d'un
facteur deux, et diviserait aussi le broadcast par deux.
